% TOP_PROBLEM: {"id": 2305037, "title": "Research Problems in Function Theory — Problem 5.37", "category_id": 8, "category": {"id": 8, "name": "analysis", "display_name": "Analysis", "description": "Limits, continuity, calculus, and function theory.", "slug": "analysis", "order_index": 8, "created_at": "2026-07-31T15:26:25.671Z"}, "status": "open"}
% FIRST_POSTED: null
% ENTRY_KIND: statement_only
% Imported from: batch14.tex; UnsolvedMath ID 2305037
% Compile twice: lualatex 2305037.tex
\documentclass[11pt]{article}
\usepackage{fontspec}
\setmainfont{Latin Modern Roman}
\usepackage{amsmath,amssymb,mathtools,unicode-math}
\setmathfont{Latin Modern Math}
\usepackage{xurl,hyperref}
\hypersetup{hidelinks,pdftitle={UnsolvedMath Problem 2305037}}
\newcommand{\sref}[2]{\hyperlink{source:#1:#2}{[S#2]}}
\newcommand{\cataloguescope}[1]{\paragraph{Mathematical question.}#1}
\begin{document}
% UnsolvedMath ID 2305037; source catalogue code AMR-022-5037
\setcounter{section}{2305036}
\section{Research Problems in Function Theory — Problem 5.37}\label{2305037}
{\small\sffamily UnsolvedMath ID 2305037 · Analysis}
\subsection{Definitions and mathematical statement}

\paragraph{Shared notation.}$\mathbb N=\{1,2,\ldots\}$, and $\mathbb Z,\mathbb Q,\mathbb R,\mathbb C$ denote the integers, rationals, reals and complex numbers. Any other conventions are specified locally.


Write $\mathbb D=\{z\in\mathbb C:|z|<1\}$ and $\mathbb T=\partial\mathbb D$. All Taylor coefficients below are taken at $0$.Let $1\le n_0<n_1<\cdots$ be integers and let $f(z)=\sum_{k\ge0}c_kz^{n_k}$ be holomorphic on $\mathbb D$, with $\sum_{k\ge0}|c_k|=\infty$. For real $\alpha<\beta$ with $0<\beta-\alpha\le2\pi$, let $V_{\alpha,\beta}=\{re^{i\theta}:0<r<1,\ \alpha<\theta<\beta\}$ be an open angular sector.Do not assume Hadamard gaps. Put $M(r,f)=\max_{|z|=r}|f(z)|$,
\[\mu=\limsup_{r\uparrow1}\frac{\log\log M(r,f)}{-\log(1-r)},\qquad N^0(t)=\#\{k:n_k\le t\}.
\]
Suppose for some $0<\gamma<1$ that $N^0(t)=O(t^{1-\gamma})$ as $t\to\infty$ and $0<\mu<(1-\gamma)/\gamma$. The positive growth index makes the double logarithm well-defined for $r$ sufficiently near $1$.
\paragraph{Statement recorded in the supplied source.}
For every finite value $a\in\mathbb C$, must (i) $f(z)=a$ have at least one solution in $\mathbb D$; (ii) it have infinitely many distinct solutions in $\mathbb D$; and (iii) it have at least one solution in every sector $V_{\alpha,\beta}$? These are three separate assertions.Investigate in particular the sequences $n_k=\lfloor(k+1)^\eta\rfloor$ with $1<\eta<3/2$, subject to the stated counting and growth conditions. \sref{2305037}{2}
\subsection{Short English statement}
For sparse gap series in the intermediate positive-growth range, ask the same three value-distribution questions, especially for exponents growing like k to a power between one and three halves.
\subsection{Sources}
\begin{description}
\item[\hypertarget{source:2305037:1}{S1}] ULAM.AI and the UnsolvedMath Contributors, UnsolvedMath: A Curated Collection of Open Mathematics Problems, record 2305037 (AMR-022-5037). CC BY 4.0. \url{https://huggingface.co/datasets/ulamai/UnsolvedMath}.
\item[\hypertarget{source:2305037:2}{S2}] W. K. Hayman and E. F. Lingham, Research Problems in Function Theory (2018), arXiv:1809.07200, Chapter 5, Problem 5.37, its update and the preceding shared notation. \url{https://arxiv.org/abs/1809.07200}.
\end{description}
\label{end:2305037}
\end{document}
